\section{All roots at small corank in the below-cap comparison range}\label{sec:small-corank-subquadratic}

2026-11-26 \enspace Private research note

There is an absolute constant $C_*>0$ such that the following holds. Let $n\ge2$, $d\ge2$, and $1\le r\le d$. If $d\le n$ and

\begin{equation}
b=\frac{d^2}{n}\ge C_*\log(n+1),
\qquad
c=d-r\le\frac{b}{1000},
\tag{T1}
\end{equation}

then the balanced-array polynomial $P(q)=\sum_{k=0}^D a_kq^k$, $D=\lfloor nr/d\rfloor$, has exactly $D$ simple negative roots. For $d>n$ the normalized coefficient separation already gives the same conclusion for every cap. The full-rank case $c=0$ is also known. Thus consider $1\le c<d\le n$ under (T1). Increase $C_*$ as needed below. Since $b\le d$, this also forces $d$ arbitrarily large once $C_*$ is chosen large.

Let $F(q)=\sum_{k=0}^n b_kq^k$ with $b_k=\binom nk^d$, and write $-\rho_i$ for its simple negative roots ordered by increasing absolute value. Set

\begin{equation*}
j_0=\lfloor d/10\rfloor,\qquad i_0=n-j_0.
\end{equation*}

Because $D=n-\lceil nc/d\rceil$ and $nc/d\le d/1000$, we have $D\ge i_0+1$ for $d$ sufficiently large. The proof transfers the first $i_0$ full-rank roots and then finds $D-i_0$ roots by alternating signs beyond the join.

\subsection{Transferring the root disks}

Use the full-rank disk of radius $\rho_i/(4n^2)$ about $-\rho_i$ for every $1\le i\le i_0$. These disks are pairwise disjoint by the logarithmic spacing estimate. At a boundary point $z$, put $t=|z|$, and let $K$ have weights $b_k t^k/F(t)$. Given $K$, sample $d$ independent uniform $K$-subsets as the columns. Write $L=n-K$ for the hole quota and $j=n-i\ge j_0$. Then

\begin{equation}
\frac{|P(z)-F(z)|}{F(t)}
\le \Pr(\text{some row has fewer than }c\text{ holes}).
\tag{T2}
\end{equation}

The root-circle bound gives

\begin{equation}
\frac{|F(z)|}{F(t)}
\ge\exp[-C_1(1+\mathrm{Var}(L))\log(n+1)] .
\tag{T3}
\end{equation}

Here and below we use the full-rank root-circle and variance estimates in the below-cap proof. To apply the local variance near the right endpoint, reflect its argument. The law of $L$ has weights proportional to $\binom n\ell^d t^{-\ell}$, so if $m$ is a mode the same mode estimate gives $\mathbb E(L-m)^2\le C_2(1+m/d)$. The law of $K$ is a sum of independent Bernoulli variables with success parameters $t/(t+\rho_j)$, and the law of $L$ is the reflected sum. On the disk boundary the number of full-rank moduli below $t$ is $i$ or $i-1$. The discrepancy of each Bernoulli parameter from its corresponding threshold indicator is at most twice its variance term. Thus $|\mathbb E L-j|\le1+2\mathrm{Var}(L)$. Combining these bounds gives $m\le j+C_3(1+m/d)$. For large enough $d$, absorb $C_3m/d$ and conclude

\begin{equation*}
m=O(j+1),\qquad
\mathrm{Var}(L)\le C_4(1+j/d),\qquad
|\mathbb E L-j|\le C_5(1+j/d).
\end{equation*}

As $j\ge\lfloor d/10\rfloor$, the last bound is $o(j)$ uniformly. In particular $\mathbb E L\ge 0.9j$ once $d$ is large. Bernstein's inequality for the Bernoulli sum, applied to a downward deviation of at least $0.4j$, gives $\Pr(L<j/2)\le e^{-c_1j}$. \href{https://github.com/ProgramComputer/open-math-research/blob/4fb624ba937a8eecc1361a6066185cd81ab7b172/topics/balanced-arrays/dense-regime/below-cap/proof.md}{Below-cap proof and root-circle bound, pinned main revision}

Given $L\ge j/2$, the hole count in any specified row has law $\mathrm{Bin}(d,L/n)$ and mean $\lambda=dL/n\ge b j/(2d)$. For large $d$, $j/d\ge0.09$, so $c\le b/1000\le\lambda/4$. A multiplicative binomial lower-tail bound gives $\Pr(\text{row has fewer than }c\text{ holes}\mid L)\le e^{-c_2\lambda}$. Union over $n$ rows and using $b\le d$ yields

\begin{equation*}
\frac{|P(z)-F(z)|}{F(t)}
\le e^{-c_1j}+n e^{-c_3 b j/d}
\le (n+1)e^{-\kappa b j/d}.
\end{equation*}

Choose $C_*$ in (T1) large enough that this is strictly smaller than (T3) using $\mathrm{Var}(L)\le C_4(1+j/d)$ and $j/d\ge0.09$. Rouché gives one root, counted with multiplicity, in each of the first $i_0$ disks. Since each disk is conjugation-invariant and lies strictly in the left half-plane, its root is simple and negative.

\subsection{The join between the disks and the upper signs}

For any positive coefficient sequence $u_k$ let $R_k=u_k/u_{k-1}$, $\Gamma_k=R_k/R_{k+1}$, and $t_k=(R_kR_{k+1})^{-1/2}$. If the ratios $R_k$ decrease, the relative magnitudes of the two neighbors of the $k$th term at radius $t_k$ are both $\Gamma_k^{-1/2}$, and every further outward term ratio is no larger. At radii $s$ with $|\log(s/t_k)|\le0.1$, every outward ratio is at most $e^{0.1}\Gamma_k^{-1/2}$.

For $F$ at $k=i_0$, its coefficient curvature satisfies

\begin{equation*}
\log \Gamma_{i_0}(F)
=d\log[(1+1/i_0)(1+1/j_0)]
> \frac{d}{j_0+1}>6
\end{equation*}

for all sufficiently large $d$. Put $t_F=\sqrt{b_{i_0-1}/b_{i_0+1}}$. On every circle of radius $s$ with $t_F e^{-0.1}\le s\le t_F e^{0.1}$, the sum of off-center term magnitudes is bounded by $2q/(1-q)$ times the central magnitude, with $q=e^{0.1-3}<1/3$. Thus the single central monomial dominates on every such circle. In particular Rouché at the lower boundary gives exactly $i_0$ full-rank roots below radius $t_F e^{-0.1}$, and hence $\rho_{i_0}<t_F e^{-0.1}$.

At each deterministic quota $k=i_0-1,i_0+1$, the hole quota $\ell=n-k$ is at least $0.08d$ for large $d$. Its specified-row mean is at least $0.08b$, so the same binomial bound gives

\begin{equation*}
1-\frac{a_k}{b_k}\le n e^{-\kappa' b}\le(n+1)^{-2}
\end{equation*}

after increasing $C_*$. Let $t_P=\sqrt{a_{i_0-1}/a_{i_0+1}}$. This deterministic acceptance estimate gives $|\log(t_P/t_F)|=O((n+1)^{-2})$. Therefore all first $i_0$ transferred disks are strictly below radius $t_P$, since their upper modulus is at most $\rho_{i_0}(1+1/(4n^2))<t_F e^{-0.09}$ for large $n$.

At every $i_0\le k<D$, the normalized log-concavity inequality for actual coefficients gives

\begin{equation*}
\log\Gamma_k(P)\ge d\log[(1+1/k)(1+1/(n-k))]
>\frac{d}{n-k+1}\ge\frac{d}{j_0+1}>6 .
\end{equation*}

The coefficient ratios are strictly decreasing throughout the support. At each geometric node $t_k(P)$, therefore, the central monomial dominates the total magnitude of the other terms by the same geometric-series bound with $q=e^{-3}$. This proves the sign $(-1)^k$ at $-t_k(P)$. The nodes strictly increase in magnitude, start at $t_P$, and the sign sufficiently far to the left is $(-1)^D$ by the leading coefficient. The $D-i_0$ intervening intervals give distinct negative roots of modulus greater than $t_P$. Together with the $i_0$ transferred roots they exhaust the degree and are simple. \href{https://github.com/ProgramComputer/open-math-research/blob/4fb624ba937a8eecc1361a6066185cd81ab7b172/topics/balanced-arrays/orientation-normalization.md}{Normalized log-concavity, pinned main revision}

\subsection{Example parameter sequences}

\begin{itemize}
\item For $d=\lceil n^{2/3}\rceil$, all coranks $1\le c\le\lfloor n^{1/3}/1000\rfloor$ are covered eventually.
\item For a fixed sufficiently large $A$, take $d=\lceil A\sqrt{n\log(n+1)}\rceil$. This covers $1\le c\le\lfloor A^2\log(n+1)/1000\rfloor$ eventually.
\end{itemize}

Both examples have $n/d^2\to0$. The argument leaves larger corank outside condition (T1) open.

\subsection{Transfer to a corank bound in the row count}

Suppose the separate universal height theorem is available in the form $n\ge C_Hd^{5/3}$ for some absolute $C_H\ge1$. Combining it with (T1) gives an absolute constant $\kappa>0$ for which every cap with $d-r\le\kappa n^{1/5}$ has all roots real, and simple negative for $d\ge2$. In fact, if $d\le(n/C_H)^{3/5}$, apply the height theorem. If $(n/C_H)^{3/5}<d\le n$, then

\begin{equation*}
b=\frac{d^2}{n}>C_H^{-6/5}n^{1/5}.
\end{equation*}

For sufficiently large $n$ this is at least $C_*\log(n+1)$. Choose $\kappa\le C_H^{-6/5}/1000$ to get $d-r\le b/1000$. The theorem above then applies with its same integer join: the cap condition gives $\lceil n(d-r)/d\rceil\le d/1000+1<\lfloor d/10\rfloor$ for large $d$, leaving the join index internal. The disk transfer and the capped-center comparison use only this cap condition and the logarithmic lower bound on $b$. For $d>n$, use normalized coefficient separation. Full rank handles the remaining finitely many row counts after decreasing $\kappa$ so that $\kappa n^{1/5}<1$ at those counts. For $d=1$, the polynomial is $(1+q)^n$.
